of finite presentation).

Since these properties are local on @@M0@@ and @@M1@@, one may suppose them affine; let @@M2@@, @@M3@@. We then have:
@@M4@@
by EGA IV 12.3.2, where the hypothesis of finite presentation on @@M5@@ can be replaced by noetherian hypotheses. Now @@M6@@ is isomorphic to the residue field of @@M7@@, and by EGA @@M8@@ 17.2.5 we have:
@@M9@@
It then follows from EGA @@M10@@ 17.2.11 that the latter is finite if and only if @@M11@@ has finite cohomological dimension, that is, is regular.

\sgaStableTarget{sga6.e09.prop.1-6}\textbf{Proposition 1.6.} \uline{Let} @@M12@@ \uline{be a noetherian scheme}, @@M13@@ \uline{a locally free} @@M14@@\uline{-Module of finite type}, @@M15@@ \uline{the vector bundle associated to} @@M16@@, \uline{and} @@M17@@ \uline{the structural morphism of} @@M18@@. \uline{Then} @@M19@@ \uline{is an isomorphism}.

Since @@M20@@ is smooth, and has a section, namely the zero section, to prove 1.6 it suffices to show that for every @@M21@@, @@M22@@ (1.4). We are thus reduced to the case where the base scheme is a field.

In this case, @@M23@@ is a direct sum of submodules of rank @@M24@@, which allows us to decompose @@M25@@ into a sequence of fibrations of rank @@M26@@, and thus reduce to the case where @@M27@@ has rank @@M28@@, namely @@M29@@. The proposition then follows from the next lemma:
